% ============================================================================
%  SECCIÓN 3.12: EJECUCIÓN DE LA APLICACIÓN
% ============================================================================

\section{Ejecuci\'on de la aplicaci\'on}

La ejecuci\'on de la aplicaci\'on se realiza sobre un tel\'efono f\'isico con
sistema operativo Android, siguiendo el flujo de desarrollo de Expo. Para ello
se utiliza el servidor de desarrollo de Expo CLI, que compila la aplicaci\'on y
la entrega al dispositivo conectado en la misma red local \citep{expo}.

El proceso de ejecuci\'on consiste en iniciar el servidor de desarrollo con el
siguiente comando, ejecutado desde el directorio del proyecto:

\begin{lstlisting}
pnpm start
\end{lstlisting}

Una vez iniciado, Expo CLI genera un c\'odigo QR que se escanea desde la
aplicaci\'on Expo Go instalada en el tel\'efono f\'isico. La aplicaci\'on se
compila y se muestra en el dispositivo, permitiendo probar las pantallas
desarrolladas en tiempo real. En la figura se presenta el servidor de
desarrollo en la terminal.

\begin{figure}[H]
\centering
  \caption{Servidor de desarrollo de Expo en la terminal}
  \label{fig:servidor-expo}
  \vspace{6pt}
  \begin{lstlisting}[style=terminal]
$ pnpm start

> lugares-interactivos@1.0.0 start
> expo start

Metro waiting on exp://192.168.1.10:8081
Logs for your project will appear below.

> Scan the QR code above with Expo Go to open the app.
> Press a | open Android

Android: opened lugares-interactivos on R58M75TGGVA
  \end{lstlisting}
  \par\raggedright\fuenteAPA{Elaboraci\'on propia en base a la ejecuci\'on del proyecto.}\end{figure}

La l\'inea \texttt{Android: opened lugares-interactivos on R58M75TGGVA}
confirma que la aplicaci\'on fue abierta en el dispositivo f\'isico conectado.
El dispositivo f\'isico es fundamental para validar el rendimiento, la
interacci\'on y el comportamiento de la interfaz en condiciones reales de uso.

En la figura se presenta la aplicaci\'on ejecutada en el dispositivo f\'isico
utilizado para el desarrollo.

\begin{figure}[H]
\centering
  \figuraAPA{fig:dispositivo-fisico}{Aplicaci\'on \nombreApp{} ejecutada en el dispositivo f\'isico}{%
    \includegraphics[width=0.62\linewidth]{appmovil-ejecutando-fisico.jpeg}%
  }{Elaboraci\'on propia (fotograf\'ia del dispositivo f\'isico).}
\end{figure}
